\chapter*{Introducción}\label{introduccion}\label{sec:introduccion}
\thispagestyle{empty}
\markboth{Introducción}{Introducción}

Un grupo que trabaja junto el tiempo suficiente recuerda más que cualquiera de sus
miembros por separado, porque cada uno sabe quién sabe qué y a quién preguntar.
Wegner~\cite{wegner1987transactive} llamó \emph{memoria transactiva} a esta
organización, y más tarde la describió como un modelo de red de
computadoras~\cite{wegner1995computer}.
Desde los sistemas computacionales, la idea es la de una red de nodos especializados
que aprenden, con el uso, a enviarse las consultas sin un coordinador central.

Las memorias asociativas entrópicas de Pineda y Morales
\cite{pineda2021entropic,pineda2022weighted,morales2025missing} guardan lo que registran
como conteos de frecuencia, reconocen una entrada cuando sus rasgos tienen soporte en
esos conteos, recuperan muestreándolos y rechazan lo que no contienen.
Este trabajo pregunta si con ellas se puede construir un sistema de memoria
transactiva, y qué agrega la memoria cuando se la compara con decisores simples que usan
las mismas representaciones.

El objetivo es construir un sistema de ocho módulos especialistas que enrute consultas
de texto e imagen, recupere el contenido asociado en ambas direcciones y rechace lo que
no conoce, sin clasificadores discriminativos en esas decisiones, y medir su desempeño.

El sistema, EAM-TMS, tiene un módulo por cada clase del conjunto ETH-80
\cite{leibe2003eth80}.
Cada módulo guarda palabras e imágenes de su clase en una sola memoria
heteroasociativa y lleva dos directorios, uno textual y uno visual.
Las palabras entran como vectores fastText~\cite{bojanowski2017fasttext} y las imágenes
como códigos de un autocodificador ResNet18~\cite{he2016resnet}.
Un mediador reparte las consultas mientras los directorios se forman; después, los
módulos se las envían entre sí.
El sistema se evalúa con un banco de 411 consultas de texto, 40 consultas fuera de
dominio, 656 imágenes de prueba y objetos que ningún módulo vio.
Se compara con decisores simples sobre los mismos vectores, con una sola memoria para
las ocho clases y con un codificador entrenado sin la pérdida de clase.

El Capítulo~\ref{cap:pp} plantea el problema, los objetivos y los alcances.
El Capítulo~\ref{sec:marco_teorico} presenta las memorias asociativas entrópicas, su
extensión heteroasociativa y el modelo de Wegner.
El Capítulo~\ref{cap:desarrollo} describe la arquitectura y la metodología experimental, el Capítulo~\ref{cap:resultados}
reporta y discute los resultados, y el Capítulo~\ref{conclusiones} responde a la
pregunta de investigación y al objetivo general, y propone el trabajo futuro.
